\documentclass[a4paper, 12pt]{article} % Adicionado tamanho do papel e fonte maior (12pt)
\usepackage[utf8]{inputenc}
\usepackage{array}
\usepackage{booktabs}
\usepackage{rotating}
% --- Configurações de Margem ---
\usepackage[a4paper, margin=2.5cm]{geometry} % É AQUI QUE A MÁGICA ACONTECE! 
% Se o professor exigir ABNT, troque a linha acima por:
% \usepackage[a4paper, left=3cm, right=2cm, top=3cm, bottom=2cm]{geometry}
\usepackage{float}
\usepackage{amsmath}

% --- Idioma e Codificação ---
\usepackage[T1]{fontenc}
\usepackage[utf8]{inputenc} % Garante que acentos funcionem (opcional nas versões mais recentes, mas seguro)
\usepackage[portuguese]{babel} % Hifenização correta e tradução (ex: "Abstract" vira "Resumo")
\usepackage{tabularx}
% --- Pacotes Úteis para Relatórios ---
\usepackage{graphicx} % Para inserção de imagens (o que você já tinha)
\usepackage{hyperref} % Deixa o sumário e as URLs clicáveis
\usepackage{microtype} % Melhora drasticamente a justificação do texto (evita buracos brancos)
\usepackage{indentfirst} % Indenta o primeiro parágrafo de cada seção (padrão do português)
\usepackage[backend=biber, style=abnt]{biblatex} % style=abnt para ABNT, ou style=ieee
\addbibresource{referencias.bib}
\usepackage{listings}
\usepackage{xcolor}

% Configuração básica para parecer com o seu exemplo
\lstset{
  basicstyle=\ttfamily\small,
  backgroundcolor=\color{gray!10},
  frame=single,
  breaklines=true
}

% --- Dados da Capa ---
\title{\textbf{Parte 3 do EP 2}}
\author{\small\textit{Matheus Rodrigues da Silva Espalaor -- 14598615}


\newcommand{\code}[1]{\texttt{#1}}
\begin{document}

\maketitle

\section{Estado da Arte}

De acordo com \cite{2023db}, atualmente, existem mais de 20 database vector systems no mercado. Vamos focar nos mais modernos e usados atualmente no mercado e no ambiente ciêntifico.

\section{Pinecone}

A evolução recente dos \textit{Vector Database Management Systems} é marcada pela transição de índices mantidos estritamente em memória para arquiteturas nativas em nuvem e sem servidor. Neste cenário, o \textbf{Pinecone} destaca-se, projetado especificamente para suprir as demandas de latência e consistência de aplicações baseadas em LLMs e RAG.~\cite{zhan2025bigvectorbench}.

O Pinecone, diferente dos próximos que serão apresentados, é de código fechado e é um  SaaS, isto é, o usuário não se preocupa em gerenciar o sistema, a escalabilidade, nem a infraestrutura. Porém, isso vem com um alto custo monetário e o suário fica preso ao sistema de nuvem do Pinecone.

\subsection{Arquitetura Desacoplada e Armazenamento}
O sistema de armazenamento utiliza estruturas de arquivos imutáveis persistidos em serviços de armazenamento de objetos distribuídos. A execução de consultas é delegada a executores, permitindo a auto-escalabilidade horizontal e independente do throughput de leitura ($QPS$) e do volume total de vetores armazenados, mitigando os custos operacionais de infraestrutura em escala de bilhões de vetores.

\subsection{Busca Híbrida e Filtragem de Metadados em Etapa Única}
Um dos principais gargalos teóricos em \textit{Vector Database Management Systems} é a eficiência em buscas vetoriais restritas por atributos escalares (filtros de metadados). As abordagens tradicionais dividem-se em:
\begin{enumerate}
    \item \textbf{Pré-filtragem:} Reduz o espaço de busca aplicando o filtro escalar antes da navegação no grafo ANN, o que frequentemente fragmenta a topologia do grafo e degrada o retorno dos dados.
    \item \textbf{Pós-filtragem:} Executa a busca vetorial no grafo completo e descarta os resultados que não satisfazem o filtro escalar, resultando em latência desproporcional e retrabalho computacional caso a seletividade do filtro seja alta.
\end{enumerate}

O Pinecone resolve esta limitação através do mecanismo de \textbf{Filtragem em Etapa Única}, onde índices invertidos e vetores de bits de metadados são integrados diretamente ao processo de travessia do algoritmo de Vizinhos Mais Próximos Aproximados (ANN). Esse fluxo unificado garante latência determinística $O(\log N)$ mesmo em cenários com filtros altamente seletivos e combinados com representações esparsas e densas simultaneamente.

\section{ Milvus }
Diferentemente do Pinecone, o Milvus é um sistema gerenciador de código aberto. Amplamente estudado e usado no meio cietifico e empresarial.

\subsection{Desacoplamento de Atores e Ciclo de Vida dos Dados}
A inovação estrutural do Milvus reside na separação estrita entre computação, armazenamento e fluxo de mensagens de log. O sistema divide suas operações em três camadas independentes:
\begin{enumerate}
    \item \textbf{Camada de Coordenação:} Funciona como o cérebro do cluster, gerenciando a atribuição de nós, topologia da rede, balanceamento de carga e metadados do sistema através de um motor de consenso distribuído.
    \item \textbf{Camada de Execução:} Composta por nós trabalhadores especializados e \textit{stateless}. Subdivide-se em \textit{Query Nodes}, exclusivos para computação de consultas ANN, \textit{Data Nodes}, responsáveis pelo consumo de logs e persistência e \textit{Index Nodes}, dedicados à construção assíncrona de índices vetoriais pesados.
    \item \textbf{Camada de Mensageria e Armazenamento:} Utiliza corretores de mensagens, como Apache Pulsar ou Kafka, para garantir a persistência de logs de escrita (\textit{Write-Ahead Logging} - WAL) antes que os dados sejam armazenados em sistemas como Amazon S3.
\end{enumerate}

\subsubsection{Consistência Dinâmica e Segmentação de Dados Semânticos}
Diferente dos bancos relacionais ACID clássicos, o Milvus permite ao usuário parametrizar o nível de consistência em tempo de execução, variando desde a consistência forte até a consistência eventual \cite{Mivus}. Para processar o fluxo contínuo de dados, o sistema introduz os conceitos de dados \textit{Bounded}, segmentos fechados e indexados e \textit{Unbounded}, dados recém-inseridos que residem em buffers de memória.

Durante uma consulta, o Milvus executa uma busca paralela híbrida. Temos o conjunto de vetores já consolidados em índices estruturados e  os vetores em estado de \textit{stream} ainda não indexados. A operação de busca de vizinhos mais próximos é formalizada como a união dos resultados computados sobre ambas as partições de memória.

Onde o KNN executa uma busca exaustiva de força bruta via força computacional paralela, usando núcleos CUDA em GPUs, se disponíveis, garantindo que dados recém-inseridos fiquem disponíveis para consulta em milissegundos sem interromper a estabilidade dos índices globais de larga escala \cite{Mivus}.


\section{ pgvector }
Diferentemente do Milvus e do Pinecone, o pgvector não foi construido do zero como um sistema puramente vetorial. Mas sim, foi construido como uma extensão do Postgres. Dessa forma, transformou o PostgreSQL em um banco de dados que suporta vetorial e relacional.

\subsection{Mecanismos de Funcionamento e Indexação}

A extensão introduziu o tipo nativo \texttt{vector} e operadores geométricos fundamentais para calcular a distância entre um vetor de consulta $\vec{q}$ e um vetor armazenado $\vec{x}$. É utilizado calculo de distância simples já mencionados anteriormente, como: distância euclidiana e a similaridade de cossenos.

\subsection{Algoritmos de Indexação Suportados}
Para mitigar a complexidade exaustiva $\mathcal{O}(N)$, o \texttt{pgvector} implementa dois algoritmos principais de Busca por Vizinhos Próximos Aproximados (\textit{Approximate Nearest Neighbor Search - ANNS}):
\begin{enumerate}
    \item \textbf{IVFFlat:} Executa a clusterização do espaço vetorial através de $k$-means criando listas invertidas. Possui baixo consumo de memória e rápida construção, mas menor precisão se o número de partições examinadas for pequeno.
    \item \textbf{HNSW:} Constrói um grafo hierárquico estruturado de múltiplos níveis. Possui uma ótima precisão e um ótimo tempo, embora demande maior processamento e memória RAM para a construção e manutenção do índice.
\end{enumerate}

\subsection{Vantagens:}
\begin{itemize}
    \item \textbf{Consistência ACID Completa:} O vetor e os dados transacionais residem no mesmo bloco de dados, garantindo atomicidade absoluta via logs de escrita antecipada (WAL).
    \item \textbf{Filtragem Híbrida Eficiente:} O otimizador de consultas do PostgreSQL avalia estatísticas estruturadas e índices vetoriais no mesmo plano de execução, permitindo instruções SQL complexas
\end{itemize}

\subsection{Limitações Estruturais do pgvector}
Apesar de sua enorme conveniência, a arquitetura baseada em extensão impõe gargalos físicos rígidos que impedem o \texttt{pgvector} de substituir sistemas massivamente distribuídos (como o Milvus) em cenários de hiperescala.

\subsubsection{Gargalos e Restrições Técnicas}

\begin{enumerate}
    \item Limite de Dimensões: o que restringe o número de dimensões do embedding para 2.000, impedindo o uso de modelos científicos pesados de última geração sem a aplicação de técnicas de compressão.
    \item Dependência de CPU: Depende inteiramente de instruções SIMD de CPU. Sistemas como o Milvus utilizam paralelismo massivo em GPU (Nvidia CUDA), operando ordens de magnitude acima do rendimento suportado pelo PostgreSQL em grandes volumes.
    \item Acoplamento de recursos: Como o PostgreSQL opera em uma arquitetura baseada em processos estruturados para CPUs, cálculos vetoriais intensivos competem diretamente por CPU e memória RAM com transações críticas de leitura/escrita. Um pico de buscas semânticas pode exaurir os recursos e degradar a aplicação principal.
\end{enumerate}

\section{ Qdrant }
O \textbf{Qdrant} é um banco de dados de vetores projetado para buscas de vizinhos mais próximos estruturadas (\textit{Approximate Nearest Neighbor - ANN}) com suporte nativo a filtragem de metadados. Desenvolvido em Rust, seu diferencial reside na eficiência de memória e na execução de filtros em estágio único.

\subsection{Estrutura de Dados Básica}
No Qdrant, a unidade fundamental de armazenamento é um ponto. Matematicamente, um ponto $P$ é definido como uma tupla:

\begin{equation}
P = \left( id, \vec{v}, \text{payload} \right)
\end{equation}

Onde:
\begin{itemize}
    \item $id \in \mathbb{N}$ ou $\text{UUID}$ é o identificador único do ponto.
    \item $\vec{v} \in \mathbb{R}^d$ é o vetor de alta dimensão ($d$). O Qdrant também suporta múltiplos vetores nomeados por ponto: $\vec{v}_1, \vec{v}_2, \dots, \vec{v}_n$.
    \item $\text{payload}$ é um objeto JSON que armazena os metadados associados (ex: marcas, preços, strings).
\end{itemize}

\subsection{Métricas de Distância}
Para determinar a similaridade entre o vetor de busca $\vec{q}$ e um vetor armazenado $\vec{v}$, o Qdrant calcula a distância matemática baseada na métrica escolhida na criação da coleção: como similaridade de cossenos, distância euclidiana, entre outros


\subsection{O Índice de Busca: HNSW}
Para evitar a busca linear $O(N)$, o Qdrant utiliza uma variação do algoritmo \textbf{HNSW} (\textit{Hierarchical Navigable Small World}). 

O HNSW constrói uma estrutura de grafos em múltiplas camadas . 
\begin{enumerate}
    \item A \textbf{Camada Superior ($L_m$)} possui poucos nós e conexões longas para buscas rápidas e grosseiras.
    \item À medida que se desce nas camadas ($L_{m-1} \dots L_0$), o grafo se torna mais denso, permitindo o refinamento da busca até encontrar os vizinhos mais próximos na $L_0$), que contém todos os vetores.
\end{enumerate}

\subsection{Filtragem em Estágio Único}
O maior diferencial do Qdrant é como ele une a busca vetorial com os filtros de \textit{payload}. Em vez de filtrar antes ou depois da busca vetorial, o Qdrant faz a \textbf{filtragem dinâmica durante a travessia do grafo HNSW}.

Uma condição de filtro definida pelo usuário, por exemplo: \texttt{preco < 50}). Durante a busca no grafo, ao avaliar se deve avançar para um nó vizinho $V_j$, o algoritmo valida: se o payload do $V_j$ satisfaz a condição do usuário

Se o nó não satisfizer a condição, a aresta é ignorada no cálculo de proximidade vetorial, garantindo que o resultado final seja matematicamente preciso e respeite estritamente as regras de negócio sem perder performance.

\subsection{Otimização de Memória: Quantização}
Para mitigar o alto consumo de RAM, o Qdrant implementa a Quantização Escalar. Ela converte os valores de ponto flutuante de 32 bits para inteiros de 8 bits, técnica similar a empregada em modelos de Deep Learning.

Isso reduz o tamanho do índice em até $4\times$, permitindo que bilhões de dimensões sejam processadas com eficiência de cache de CPU e consumo de memória drasticamente menor.

\newpage

\section{Análise Comparativa}

A tablea apresenta uma consolidação visual e técnica dos principais sistemas gerenciadores de bancos de dados vetoriais analisados neste relatório.

\begin{table}[H]
\centering
\caption{Comparativo: Linguagem Base e Algoritmos de Indexação}
\label{tab:tecnologia_indexacao}
\small
\setlength{\extrarowheight}{3pt}
\begin{tabularx}{\textwidth}{l l X}
\toprule
\textbf{Sistema} & \textbf{Linguagem Base} & \textbf{Algoritmos de Indexação Suportados} \\
\midrule
\textbf{Pinecone} & Go / C++ & Proprietário (Estrutura de Grafo / ANN). \\
\addlinespace
\textbf{Milvus}   & Go / C++ / Python & HNSW, IVF, ScaNN, ANNOY. \\
\addlinespace
\textbf{pgvector} & C (Extensão Postgres) & HNSW, IVFFlat. \\
\addlinespace
\textbf{Qdrant}   & Rust & HNSW customizado. \\
\bottomrule
\end{tabularx}
\end{table}

% --- TABELA 2: BUSCA HÍBRIDA E LICENÇA ---
\begin{table}[H]
\centering
\caption{Comparativo: Suporte a Busca Híbrida e Tipo de Licença}
\label{tab:busca_licenca}
\small
\setlength{\extrarowheight}{3pt}
\begin{tabularx}{\textwidth}{l X l}
\toprule
\textbf{Sistema} & \textbf{Busca Híbrida (Vetor + Metadados)} & \textbf{Tipo de Licença} \\
\midrule
\textbf{Pinecone} & Filtragem em Etapa Única nativa com vetores de bits. & Proprietária (SaaS) \\
\addlinespace
\textbf{Milvus}   & Filtragem dinâmica paralela sobre dados indexados e em \textit{stream}. & Open Source\\
\addlinespace
\textbf{pgvector} & Filtros relacionais e vetoriais integrados no otimizador SQL. & Open Source\\
\addlinespace
\textbf{Qdrant}   & Filtragem dinâmica em estágio único durante a travessia do grafo. & Open Source\\
\bottomrule
\end{tabularx}
\end{table}

% --- TABELA 3: ESCALABILIDADE ---
\begin{table}[H]
\centering
\caption{Comparativo: Modelo de Escalabilidade}
\label{tab:escalabilidade}
\small
\setlength{\extrarowheight}{4pt}
\begin{tabularx}{\textwidth}{l X}
\toprule
\textbf{Sistema} & \textbf{Modelo de Escalabilidade} \\
\midrule
\textbf{Pinecone} & Nuvem nativa, totalmente desacoplada e sem servidor (\textit{Serverless}), focada em auto-escalabilidade horizontal. \\
\addlinespace
\textbf{Milvus}   & Altamente distribuído com separação estrita de três camadas independentes (Coordenação, Execução e Mensageria). \\
\addlinespace
\textbf{pgvector} & Essencialmente vertical, operando acoplado diretamente aos recursos físicos e processos de CPU/RAM da instância PostgreSQL. \\
\addlinespace
\textbf{Qdrant}   & Escalabilidade horizontal nativa através de um sistema robusto e eficiente de \textit{sharding} construído em Rust. \\
\bottomrule
\end{tabularx}
\end{table}

\begin{table}[H]
\centering
\caption{Comparativo: Vantagens e Desvantagens dos Sistemas}
\label{tab:vantagens_desvantagens}
\small
\setlength{\extrarowheight}{4pt}
\begin{tabularx}{\textwidth}{l X X}
\toprule
\textbf{Sistema} & \textbf{Vantagens} & \textbf{Desvantagens} \\
\midrule
\textbf{Pinecone} & 
Totalmente gerenciado (SaaS); infraestrutura e escalabilidade transparentes; excelente para RAG com filtragem em etapa única nativa. & 
Código fechado (proprietário); alto custo financeiro em larga escala; dependência total do ecossistema de nuvem do provedor (\textit{vendor lock-in}). \\
\addlinespace \hline \addlinespace
\textbf{Milvus} & 
Código aberto; altamente distribuído e robusto; permite consistência dinâmica e paralelismo massivo utilizando hardware focado (GPUs/CUDA). & 
Arquitetura complexa com muitas partes móveis (três camadas distintas); curva de aprendizado e manutenção de infraestrutura elevadas. \\
\addlinespace \hline \addlinespace
\textbf{pgvector} & 
Consistência ACID completa; integração nativa com dados relacionais existentes no PostgreSQL; sem necessidade de gerenciar um novo banco de dados. & 
Limite físico de 2.000 dimensões; dependência estrita de CPU (sem aceleração por GPU); cálculos vetoriais intensivos competem por recursos com transações SQL críticas. \\
\addlinespace \hline \addlinespace
\textbf{Qdrant} & 
Desenvolvido em Rust (alta eficiência de memória); filtragem dinâmica em estágio único durante a travessia do grafo HNSW; quantização escalar nativa. & 
Comunidade e ecossistema de ferramentas integradas ligeiramente menores quando comparados ao Milvus ou soluções de mercado mais antigas. \\
\bottomrule
\end{tabularx}
\end{table}

\section{Referências}
\printbibliography

\end{document}
